\documentclass[11pt]{article}
\usepackage[margin=1in]{geometry}
\usepackage{amsmath,amssymb,amsthm}
\usepackage[hidelinks]{hyperref}

\newtheorem{theorem}{Theorem}
\newtheorem{proposition}[theorem]{Proposition}
\newtheorem{lemma}[theorem]{Lemma}
\newtheorem{corollary}[theorem]{Corollary}
\theoremstyle{definition}
\newtheorem{remark}[theorem]{Remark}

\newcommand{\R}{\mathbb{R}}
\newcommand{\norm}[1]{\left\|#1\right\|}
\newcommand{\abs}[1]{\left|#1\right|}
\renewcommand{\d}{\,d}

\title{Sharp Faber--Krahn stability via the Kohler--Jobin transfer}
\author{Plan 1 / Sharp stability for Faber--Krahn}
\date{2026-05-11}

\begin{document}
\maketitle

\begin{abstract}
The spectral closure of \texttt{fixed-domain-bernoulli-expansion.md} gives
sharp \emph{Saint--Venant} stability:
\[
E(\Omega)-E(B_1)\ge c_*(N,R)\,\alpha(\Omega)
\quad\hbox{for }Q_\alpha(\Omega)\le q_*(N,R).
\]
By BDV Lemma 4.2 ($\mathcal A(\Omega)^2\lesssim\alpha(\Omega)$) and the small-vs-large
asymmetry split, this gives the global sharp Saint--Venant inequality
$\mathcal A(\Omega)^2\le C_{\rm SV}(N,R)(E(\Omega)-E(B_1))$.

This note records the Kohler--Jobin transfer that converts this into sharp
\emph{Faber--Krahn} stability for the first Dirichlet eigenvalue:
\[
\lambda_1(\Omega)-\lambda_1(B^*)\ge c_{\rm FK}(N,R)\,\mathcal A(\Omega)^2,
\]
where $B^*$ is the ball with $\abs{B^*}=\abs{\Omega}$. The transfer uses
only the Kohler--Jobin inequality (a non-perturbative pointwise comparison
that holds for every open set) plus elementary scaling. The output is the
BDV sharp Faber--Krahn theorem.
\end{abstract}

\section{The Kohler--Jobin inequality}\label{sec:kj}

For an open bounded set $\Omega\subset\R^N$, let
\[
T(\Omega):=\int_\Omega u_\Omega\d x,\qquad
-\Delta u_\Omega=1\hbox{ in }\Omega,\;u_\Omega=0\hbox{ on }\partial\Omega,
\]
be the torsional rigidity, and let $\lambda_1(\Omega)$ be the first Dirichlet
eigenvalue of $-\Delta$ on $\Omega$. Both functionals scale as
\[
T(t\Omega)=t^{N+2}T(\Omega),\qquad\lambda_1(t\Omega)=t^{-2}\lambda_1(\Omega),
\]
so the product
\[
\Psi(\Omega):=T(\Omega)\cdot\lambda_1(\Omega)^{(N+2)/2}
\]
is scale-invariant.

\begin{theorem}[Kohler--Jobin, 1978]\label{thm:kj}
For every open bounded $\Omega\subset\R^N$ with $\lambda_1(\Omega)<\infty$,
\[
\Psi(\Omega)\ge\Psi(B),
\]
where $B$ is any ball, with equality iff $\Omega$ is a ball (up to translation
and null sets).
\end{theorem}

See Kohler--Jobin~(1978) and Brasco, ``On torsional rigidity and principal
frequencies'' (ESAIM:COCV, 2014) for proofs via rearrangement. The constant
$\Psi(B)$ is universal:
\[
\Psi(B)=T(B)\cdot\lambda_1(B)^{(N+2)/2}=c_N
\]
depends only on $N$ (not on the radius), since $\Psi$ is scale-invariant.

\section{Transfer at fixed volume}\label{sec:transfer}

Fix $\abs{\Omega}=\abs{B^*}$ where $B^*$ is the comparison ball. Setting
$L:=\lambda_1(B^*)$, $T_0:=T(B^*)$, and abbreviating $\beta:=(N+2)/2$,
Theorem~\ref{thm:kj} reads
\begin{equation}\label{eq:kjfixed}
T(\Omega)\cdot\lambda_1(\Omega)^\beta\ge T_0\cdot L^\beta.
\end{equation}

\begin{lemma}[Pointwise transfer]\label{lem:pointwise}
For every $\Omega$ with $\abs{\Omega}=\abs{B^*}$,
\[
\lambda_1(\Omega)\ge L\cdot\Bigl(\frac{T_0}{T(\Omega)}\Bigr)^{1/\beta}.
\]
Equivalently, with the Saint--Venant deficit
$\delta_{\rm SV}(\Omega):=T_0-T(\Omega)\ge0$ (note $\delta_{\rm SV}\ge0$ at fixed volume by
Saint--Venant),
\[
\lambda_1(\Omega)-L\ge L\cdot\Bigl[\Bigl(\frac{1}{1-\delta_{\rm SV}/T_0}\Bigr)^{1/\beta}-1\Bigr].
\]
\end{lemma}

\begin{proof}
Rearrange~\eqref{eq:kjfixed}.
\end{proof}

\begin{lemma}[Small-deficit linearization]\label{lem:linear}
For $\delta_{\rm SV}\le T_0/2$,
\[
\lambda_1(\Omega)-L\ge\frac{L}{\beta T_0}\cdot\delta_{\rm SV}
=\frac{2L}{(N+2)T_0}\cdot\bigl(T_0-T(\Omega)\bigr).
\]
\end{lemma}

\begin{proof}
For $t\in[0,1/2]$,
$(1-t)^{-1/\beta}-1\ge t/\beta$ by elementary convexity. Apply with
$t=\delta_{\rm SV}/T_0$.
\end{proof}

\begin{remark}
The factor $2L/((N+2)T_0)$ is universal in $N$ once normalized. For the unit
ball $B_1$ in $\R^N$, $L=L_N$ is the first Dirichlet eigenvalue of $-\Delta$
on $B_1$ and $T_0=T_N$ is the torsional rigidity of $B_1$. Both are explicit:
$L_N$ is the smallest positive zero of a Bessel function, $T_N=\abs{B_1}/(N(N+2))$.
\end{remark}

\section{Sharp Faber--Krahn from sharp Saint--Venant}

\begin{proposition}[Small-asymmetry Faber--Krahn]\label{prop:smalla}
Assume the sharp Saint--Venant stability holds in the form
\begin{equation}\label{eq:svsharp}
T(B^*)-T(\Omega)\ge c_{\rm SV}(N,R)\,\mathcal A(\Omega)^2
\quad\hbox{for all }\Omega\subset B_R\hbox{ with }\abs{\Omega}=\abs{B^*}\hbox{ and }\mathcal A(\Omega)\le\mathcal A_*.
\end{equation}
Then for the same $\Omega$,
\[
\lambda_1(\Omega)-\lambda_1(B^*)
\ge\frac{2L\,c_{\rm SV}(N,R)}{(N+2)T_0}\,\mathcal A(\Omega)^2.
\]
\end{proposition}

\begin{proof}
Combine Lemma~\ref{lem:linear} with~\eqref{eq:svsharp}. The condition
$\delta_{\rm SV}\le T_0/2$ in Lemma~\ref{lem:linear} is satisfied for
$\mathcal A_*$ small enough.
\end{proof}

\begin{theorem}[Sharp Faber--Krahn stability]\label{thm:fk}
There is $c_{\rm FK}(N,R)>0$ such that, for every open $\Omega\subset B_R$
with $\abs{\Omega}=\abs{B^*}$,
\[
\lambda_1(\Omega)-\lambda_1(B^*)\ge c_{\rm FK}(N,R)\,\mathcal A(\Omega)^2.
\]
\end{theorem}

\begin{proof}
\emph{Small asymmetry.} For $\mathcal A(\Omega)\le\mathcal A_*(N,R)$,
combine our sharp Saint--Venant stability (the global form derived in
\texttt{fixed-domain-bernoulli-expansion.md}, Theorem~6.3, combined with
the small-vs-large split using BDV Lemma 4.2) with Proposition~\ref{prop:smalla}.
This gives $\lambda_1(\Omega)-L\ge c_1(N,R)\mathcal A(\Omega)^2$ with
$c_1=2Lc_{\rm SV}/((N+2)T_0)$.

\emph{Large asymmetry.} For $\mathcal A(\Omega)\ge\mathcal A_*$, the
classical Faber--Krahn inequality $\lambda_1(\Omega)\ge L$ with strict
inequality away from the ball combined with the qualitative continuity of
$\lambda_1$ in Hausdorff convergence (or any compactness argument) yields a
uniform lower bound
\[
\lambda_1(\Omega)-L\ge\inf_{\mathcal A_*\le\mathcal A(\Omega)\le 2}
\bigl(\lambda_1(\Omega)-L\bigr)=:\delta_*(N,R,\mathcal A_*)>0
\ge\delta_*\cdot\frac{\mathcal A(\Omega)^2}{4}.
\]
(Here $\mathcal A\le 2$ since $\abs{\Omega\Delta B^*}\le 2\abs{B^*}$, so
$\mathcal A^2\le 4$.) Setting $c_2:=\delta_*/4$,
\[
\lambda_1(\Omega)-L\ge c_2\mathcal A(\Omega)^2.
\]

Take $c_{\rm FK}:=\min(c_1,c_2)$.
\end{proof}

\section{Quantitative dependence on $(N,R)$}

The constant $c_{\rm FK}(N,R)$ is the minimum of the small-asymmetry and
large-asymmetry constants:

\begin{itemize}
\item \textbf{Small-asymmetry constant:}
\[
c_1(N,R)=\frac{2L_N\,c_{\rm SV}(N,R)}{(N+2)T_N},
\]
where $c_{\rm SV}(N,R)$ is the Saint--Venant constant from our spectral
closure. Tracing back through the selection chain:
\begin{itemize}
\item Selection (\texttt{quantitative-selection-principle.md}) preserves
the quotient $Q_\alpha$ up to a constant $C_{\rm sel}(N,R)$, gaining
$\sigma\le C\tau=Cq_0^{1/4}$.
\item Graph entry (\texttt{graph-entry-closure.md}) requires
$\alpha(\widetilde\Omega)\le\alpha_{\rm graph}(N,R)$.
\item Spectral closure (\texttt{fixed-domain-bernoulli-expansion.md}) requires
$\sigma\le\sigma_*(N,R)$ and $\norm{g}_{C^{2,\gamma}}\le\delta_*(N,R)$.
\end{itemize}
Combining, $q_*\sim\sigma_*^4$ (the binding constraint), and $c_{\rm SV}\ge q_*/C_{\rm sel}$.
Hence
\[
c_1(N,R)\ge \frac{2L_N\sigma_*(N,R)^4}{(N+2)T_N\cdot C_{\rm sel}(N,R)}.
\]
\item \textbf{Large-asymmetry constant} $c_2(N,R)$: a qualitative
compactness constant. Explicit lower bound (in $N=2,3$) is obtained from
quantitative versions of Faber--Krahn for the convex hull or convex
envelope, but for our purposes any positive lower bound suffices.
\end{itemize}

\begin{remark}[Comparison with BDV]
This is exactly the chain BDV use to prove sharp Faber--Krahn: their Theorem
1.1 is our Theorem~\ref{thm:fk}. The new content of Plan 1 is that
$c_{\rm SV}(N,R)$ is now derived from the Bernoulli spectral closure
(\texttt{fixed-domain-bernoulli-expansion.md}) rather than from BDV's nearly
spherical second variation. The Kohler--Jobin transfer step is shared with
BDV and uses no new input.
\end{remark}

\section{Status}

The full chain for sharp Faber--Krahn stability is structurally complete,
with two possible closure routes:
\begin{enumerate}
\item \textbf{Route A (unconditional):} BDV's nearly spherical second
variation supplies the sharp Saint--Venant constant $c_{\rm SV}(N,R)$.
\item \textbf{Route B (conditional):} the Bernoulli spectral closure of
\texttt{fixed-domain-bernoulli-expansion.md}, supplemented by
\texttt{corrections-response.md}, supplies $c_{\rm SV}(N,R)$, conditional on
the explicit constants in the second-variation source enumeration.
\end{enumerate}
Either way, the Kohler--Jobin transfer of Lemma~\ref{lem:pointwise} +
Lemma~\ref{lem:linear} + small/large split (Theorem~\ref{thm:fk}) converts
$c_{\rm SV}$ into $c_{\rm FK}$.

The quantitative Faber--Krahn constant $c_{\rm FK}(N,R)$ is given by an
explicit formula in terms of:
\begin{itemize}
\item the universal ball constants $L_N=\lambda_1(B_1)$ and
$T_N=T(B_1)=\abs{B_1}/(N(N+2))$;
\item the BDV regularity constants from selection and graph entry;
\item the spectral-closure smallness constants $\sigma_*,\delta_*$.
\end{itemize}

The exponent is sharp: the inequality
$\lambda_1(\Omega)-\lambda_1(B^*)\ge c_{\rm FK}\mathcal A(\Omega)^2$ is the
optimal power of $\mathcal A$, as already observed by Bhattacharya, Hansen,
Nadirashvili (and saturated by ellipsoidal perturbations of the ball).

\end{document}
